\section{Methods}
\label{sec:methods}

\begin{figure}[htbp]
  \centering
  \includegraphics[width=\linewidth]{P1-1_protocol_diagram.pdf}
  \caption{Protocol diagram: host custody (fingerprint-bound cases and gold labels) $\rightarrow$ frozen case $\rightarrow$ system under test (public view only) $\rightarrow$ scoring pipeline. Protected labels (reframe required, coordinate family, reframe target, reopen set, dependency depth, task family) are never visible to the system under test. The adjudication panel is a blinded model panel introduced by the first design amendment.}
  \label{fig:P1-1}
\end{figure}

This section describes the executed design. Every quantity below is fixed by
artifacts committed before any outcome was inspected; the protocol records that
no outcome had been accessed at the freeze, and the dataset is bound by content
hash, so the design cannot be reconstructed after the fact to fit a result.
Machine-readable artifacts govern rather than this narrative, and
Section~\ref{sec:availability} names where each is archived.

\subsection{Three-valued outcomes}

Every decision reported in this paper is three-valued rather than binary. A case
may be determined to have recovered, determined to have failed, or left
\emph{undetermined}; a gate may be determined to hold, determined to be
violated, or left undetermined. ``We could not measure this'' is recorded,
reported and counted separately from ``we measured this and it is false'', and
the two are never merged. The distinction is load-bearing rather than cosmetic:
an unreachable provider, an absent credential or a comparison in which every arm
produced the same outcome vector yields an undetermined result, and an
undetermined result is not evidence of absence. Wherever this paper describes an
outcome as \emph{undetermined}, that is this third value, carried in the evidence
record and in every downstream count, never silently resolved to a negative.

\subsection{Frozen protocol and bound artifacts}

The design was frozen and versioned before execution, and
Section~\ref{sec:availability} names the freeze and its amendments. The first
amendment removed human adjudication as a requirement and installed a blinded
model panel as the adjudication standard; it was admitted while no outcome had
been accessed, so it changed the design before any outcome was visible rather
than after. The amendment did not relax provenance: panel output is recorded as
model-derived adjudication, which the record type mechanically refuses to report
as human adjudication, and adjudicated labels are described as model-derived
throughout.

The evaluation suite is bound by content hash over cases and gold, separately for
the pilot and test splits. A result reported against an unbound dataset is
unverifiable, so the hashes are carried in the frozen execution bindings rather
than in prose.

\subsection{Task suite}

Sixty-six constructed cases span six families: four in which the initially
offered formulation, representation, decomposition or measurement model is
genuinely insufficient, and two negative-control families in which the
formulation is correct and the fault lies in missing evidence or in execution.
\PilotCaseCount{} cases form a pilot split used only for debugging and variance
estimation; \TestCaseCount{} form the frozen test split. Gold dependency depth
ranges over one to three so that reopening can be analysed by depth.

Each case exposes a public view: a prompt, a set of observable resources, and a
budget class. Six labels are withheld in host custody---whether a reframe is
required, the responsible coordinate family, the reframe target, the set of
closures that must reopen, the dependency depth, and the task family itself. The
family is withheld because four of six families are hidden-shift and two are
controls, so naming the family restates the answer to the restraint hypothesis.
Withholding is enforced by the type signature rather than by convention: a system
under test receives the public view and never an object carrying gold, and a
single scoring module is the only code permitted to read it.

\subsection{Suite validity}

The preregistered validity battery was applied to the then-frozen 66-case
historical corpus. Sixteen derived surface statistics were tested under two
statistics (maximum between-family mean gap; leave-one-out
nearest-class-mean accuracy) and two labellings (six-family; control versus
hidden) on both splits, against shuffle-equal-$n$ nulls; no comparison survived
Holm correction. A joint-space leave-one-out nearest-neighbour classifier over
the $z$-scored derived-feature space recovered the family at $0.062$ against a
$0.167$ majority baseline. The registered token-presence scan over the
then-frozen public prompts and observable-resource descriptions also reported
no separator, and a degeneracy probe over case features and blind responders
returned clean on both splits.

That negative surface report did not establish a shortcut-free suite. A later
author-lane solvability audit examined the historical \texttt{PublicView}
representation available to a blind responder, including exact prompt phrases
and resource-path stems. It found that ``The team's framing is'' separated all
44 hidden-shift cases from all 22 controls, while \texttt{proposal}/
\texttt{trial} path stems identified 10 of 11 decomposition cases and none of
the other 55. A responder using only path stems and two fixed phrases could
therefore solve 33 of 66 cases. The preregistered surface audit had missed these
later-audited perfect shortcuts; its negative finding is retained as chronology,
not as evidence that the historical public views were cue-free.

Before the freeze, two other exploitable rules were found and removed during
construction: a resource-count imbalance arising from balancing two members of
a linearly dependent triple, and a rule keyed on the presence of a digit in the
prompt that separated hidden cases from controls at $0.771$ against a $0.667$
baseline. The earlier audit also reported residual low-level lexical cues at
approximately $0.10$ lift. Those repairs and residual-cue estimates did not
remove the phrase and path-stem shortcuts found later. The 66-case corpus is
therefore retained as a development falsifier, not promoted as a clean
predictive benchmark or evidence of general responsibility inference.

\subsection{Systems under test}

The implementation evaluated here is the ORION reconstruction layer; it is the
artifact under test, and the rest of the paper calls it the \emph{governed
policy} so that each claim reads as a claim about the mechanism rather than
about a product.

Twelve systems are compared. Five are protocol-matched reimplementations of
published designs, constructed where original code was unavailable: a static
tool workflow, an iterative tree-search workflow, a recursive audit-and-follow-up
controller in the style of AREX \cite{arex2026}, a dependency-aware execution
plan in the style of SCION \cite{scion2026}, and a revisable
information-state controller in the style of Iris \cite{iris2026}. The governed
policy is compared against five ablations of itself: without explicit $W$,
without explicit $M$, generic retry in place of typed responsibility-targeted
reframing, full reset in place of dependency-directed reopening, and no mechanic
self-audit. An additional baseline answers from the same live provider
(glm-5.2) in a single shot with the governing control architecture removed,
providing a fresh-provider reference from the same model family.

All systems share one perception layer, so no system can win by perceiving
better rather than deciding better, and all are resource-matched, with the
single exception of the live-provider baseline, which receives the same prompt
and resource listing once without tool access. Root success is never
self-asserted: each system re-runs the shared detector against its own
post-repair state, and a formulation fault clears only if the policy actually
moved the coordinate responsible for it.

\subsection{Statistical analysis}

The analysis unit is the frozen case, with five stochastic repetitions nested
within case and system and reduced before any interval is computed; $n$ is
therefore the number of cases and never cases~$\times$~seeds. Standalone binary
rates receive Wilson score intervals. Matched differences use a paired
percentile bootstrap with ten thousand resamples over the frozen case pairing.
Effect sizes accompany every comparison. The primary hypothesis carries a
prospective superiority margin of $+\POneHOneMargin{}$ absolute root success;
the non-inferiority margin for unnecessary reframing is $+\POneHTwoMargin{}$.
Secondary inferential families are Holm-corrected. An interval that includes its
margin does not report support.

No candidate-caused failure, timeout, malformed action, abstention or incorrect
reframe is excluded. Abstention is a distinct outcome that partitions with
reframe and restraint, so a system cannot obtain a clean unnecessary-reframe
rate by declining to answer. A case that could not be measured is carried as
undetermined and leaves both numerator and denominator; it is never aggregated
as a zero.

Two admissibility conditions are evaluated before any system comparison is
reported. First, non-compensatory blocking gates: a failed or unobserved
requirement cannot be offset by performance elsewhere. Second, arm
discrimination: if every compared system produced an identical outcome vector,
the comparison reports a difference of exactly zero by construction rather than
by measurement, and the honest verdict is that the comparison is undetermined
rather than null. This condition is checked and reported, not assumed.

\subsection{Reference levels}

Two null levels are computed from the suite rather than adopted by convention.
For reopening, the mechanism-free reference is a blind policy that reopens the
largest closure component without reading the prompt or naming a
responsibility; it scores $0.792$ on pilot and $0.823$ on test, above the
full-reset ablation at $0.719$. A reopen result between those levels reflects
graph shape rather than dependency-directed reasoning, so the ablation alone is
not a sufficient comparison. For attribution, blind responders reading only case
identifiers, surface counts, or structural shape are reported alongside the
majority baseline. The frozen archive additionally records
\StochasticRepeats{} stochastic repeats per case per system; intervals reported
in Table~\ref{tab:P1-T2} are taken on the reduced case unit with Wilson score
intervals at the protocol confidence, and after the live-provider recovery the
archive carries no undetermined cell, the 86 cells lost to provider failures
having been re-executed against the frozen subject.
